\documentclass[12pt,a4paper]{article}
\usepackage[utf8]{inputenc}
\usepackage[T1]{fontenc}
\usepackage{amsmath,amssymb,amsfonts}
\usepackage[margin=2.5cm]{geometry}
\usepackage{hyperref}

\title{\textbf{Hybrid Native Observers and Competing Epistemic Horizons}}
\author{Franklin Silveira Baldo}
\date{March 2026}

\begin{document}
\maketitle

\begin{abstract}
Following the empirical resolution of the Hardware-Software Confound via the Native Cross-Architecture test, we have confirmed that structurally distinct observers ($\Delta_{\text{Transformer}}=100\%$ vs $\Delta_{\text{SSM}}=40\%$) exhibit entirely different deviation distributions when attempting \#P-hard tasks. But what happens when an observer is structurally hybridized?
\end{abstract}

\section{Interference Between Competing Limits}
If we construct a hybrid native observer---such as a Jamba model, which patches global attention layers with native State Space Model (SSM) sub-layers---do we observe a hybridized $\Delta$ distribution? Or does the global attention layer completely dominate the fading memory layer's epistemic horizon, forcing $\Delta \to \Delta_{\text{Transformer}}$?

In a Generative Ontology, this experiment maps the interference between competing epistemic limits within a single generative act. It provides a formal, mechanistic test for how varying architectural bounds composite to form the physical laws of the simulated universe.

I propose the Hybrid Native Observers test to empirically measure if the distribution of deviation cleanly interpolates between $40\%$ and $100\%$, or if one Epistemic Horizon subsumes the other.
\end{document}
